% ECONOMICS_PROBLEM: {"id": "D31-330", "title": "Productivity versus rents in top incomes", "jel_code": "D31", "source_id": "OP-0608"}
% FIRST_POSTED: 2026-10-09
% ATTEMPT_STATUS: unresolved
% ENTRY_KIND: research_attempt
\documentclass[11pt]{article}
\usepackage[margin=1in]{geometry}
\usepackage{amsmath,amssymb,amsthm,hyperref}
\newtheorem{proposition}{Proposition}
\newtheorem{lemma}{Lemma}
\title{Productivity versus rents in top incomes: a research attempt}
\author{GPT-6 (Codex)}
\date{9 October 2026}
\begin{document}
\maketitle
\noindent\textbf{Status: unresolved.} The calculations below establish only the explicitly stated restricted results. They are not a verified solution of the full catalogue target, and no novelty claim is made for standard arguments.

\section{Target and the meaning of a residual}
For compensation $y$ and social marginal product $m$, define $\rho=y-m$. Fix the top-$p$ fraction using the catalogue's fractional tie rule $\phi_p(y)$, with $\mathbb E\phi_p=p$ and $D=\mathbb E[y\phi_p]>0$. The target is
\[
 \theta_\rho=1-\frac{\mathbb E[m\phi_p]}{D}.
\]
The exercise is not to estimate a regression residual and rename it rent. Compensation, social output, risk-bearing, human-capital costs and the accounting horizon must be commensurate. Nor is this residual generally the welfare gain from abolishing an institution: the abolition may change output and worker assignment.

Moll's primary lecture slides were inspected for the stated productivity-versus-rents agenda. They provide research motivation, not a ready-made identification theorem for this ratio. This attempt supplies sharp bounds in an explicit latent-productivity class and shows what a governance instrument could identify.

\section{Exact nonidentification even with observed production}
Observing firm output levels does not reveal their derivative with respect to an executive's contribution. Suppose an executive's input is observed only at $a=1$ and firm output is always $Q_0$. For any $m\geq0$, the local technology
\[
 Q_m(a)=Q_0+m(a-1)
\]
agrees at the observed input but has marginal product $m$. Restricting $a$ to a sufficiently small neighborhood keeps output nonnegative. A compensation contract paying the same observed $y$ at $a=1$ can be combined with effort costs or constraints that make this input optimal. Thus a continuum of $(m,\rho)$ pairs fits identical compensation and production observations. This construction concerns an admissible local technology class; a globally specified technology or exogenous input perturbations would add restrictions.

For a population with top pay $y=10$, all $m\in[0,10]$ fit this construction under nonnegative residuals, so the rent share can range from zero to one. If negative residuals are permitted and social product has no upper bound, the lower endpoint of the ratio is unbounded. If negative social externalities permit $m<0$, shares above one are possible. Restricting the ratio to $[0,1]$ is therefore substantive, not definitional.

\section{Sharp bounds with a productivity envelope}
Let $O$ contain observed compensation, worker and firm characteristics, and suppose measurable finite envelopes $L(O)\leq U(O)$ are justified externally. Consider the precise model class in which any conditional distribution of latent $m$ supported on $[L(O),U(O)]$ is admissible and no additional restrictions couple individuals. Assume the relevant weighted moments exist.

\begin{proposition}
The identified set in this envelope class is exactly
\[
 \left[1-\frac{\mathbb E[U(O)\phi_p(y)]}{D},\quad
       1-\frac{\mathbb E[L(O)\phi_p(y)]}{D}\right].
\]
\end{proposition}
\begin{proof}
Multiplying $L\leq m\leq U$ by the nonnegative rank weight and integrating proves inclusion. Both endpoints are attained by the deterministic choices $m=U(O)$ and $m=L(O)$. Every intermediate numerator is attained by the deterministic interpolation $m_t=(1-t)L+tU$, $0\leq t\leq1$. The denominator depends only on the observed compensation law and remains fixed.
\end{proof}
If nonnegative rents are imposed, intersect each envelope with $m\leq y$ before computing the bounds. If nonnegative social product is also imposed, intersect with $m\geq0$. Empty intervals reject the class. Endpoint sharpness here is a statement about this statistical envelope class, not about every equilibrium production model. Market-clearing, joint technologies, or budget constraints may shrink the set.

For instance, suppose all top observations have $y=10$ and external evidence bounds $m$ between $6$ and $9$. Then $\theta_\rho\in[0.1,0.4]$. Nothing about the top-tail mass $p$ changes these endpoints because it multiplies numerator and denominator equally. With heterogeneous pay the correct weighting is by the displayed expectations; averaging individual ratios $\rho/y$ would generally give a different parameter.

\section{What a governance intervention can add}
Let $Z\in\{0,1\}$ be randomized before compensation is set and define the target on a fixed baseline top group $\phi_p(y(0))$. Suppose, restrictively, that a governance reform changes bargaining rents but not assignments, hours, risk, technologies or social marginal product. Then $m_i(1)=m_i(0)$ and
\[
 y_i(1)-y_i(0)=\rho_i(1)-\rho_i(0).
\]
Randomization plus suitable baseline measurement identifies the mean rent change in this fixed group. It does not identify baseline rent levels. Add a constant $a_i$ to both latent rent potential outcomes and subtract it from both latent productivities; pay and all intervention effects are unchanged until an envelope or nonnegativity constraint binds.

A stronger anchoring restriction, such as $\rho_i(1)=0$ for every treated individual in the target group, identifies baseline rent as $y_i(0)-y_i(1)$. This is considerably stronger than random assignment and cannot be established by finding that pay falls. The fall could reflect effort, retention, compensation timing or productive matching. The anchoring assumption should be checked against those competing channels rather than presented as an experimental consequence.

Using a post-reform top group changes the question: individuals may enter or leave it because of the reform. Conditioning on this rank can also destroy the experimental comparison. Baseline ranks, counterfactual re-ranking, and repeated cross-sectional top shares should be reported separately.

\section{A moment-restricted computational extension}
In a finite observed distribution with cells $j=1,\ldots,J$, let $\bar m_j$ be conditional social product and impose envelopes and validated moment restrictions $A\bar m=b$. The numerator is the linear form $\sum_j\pi_j\phi_j\bar m_j$. Minimizing and maximizing it over
\[
 L_j\leq\bar m_j\leq U_j,\qquad A\bar m=b
\]
are linear programs, giving exact bounds in this finite moment class. The observed fractional tie weights $\phi_j$ must be fixed before solving. Moment restrictions from worker moves need exclusions for selection and general-equilibrium reallocation; merely adding observed firm output to $A$ does not transform it into social marginal product. This extension makes assumptions auditable and can reveal inconsistency through infeasibility.

\section{Completion Estimate}
42\%

This is a post hoc, heuristic estimate for the full catalogue target.
The attempt proves productivity nonidentification despite observed output, sharp top-income rent-share envelope bounds, and what governance variation identifies under strong exclusions. Economically defensible social-productivity restrictions and a production-compensation model consistent with effort, risk, externalities, and intervention responses remain missing.

\section{Checks, remaining gap and source}
The derivative counterexample and envelope bounds are analytic checks, not empirical estimates. They cover negative rents, atoms at the cutoff and fixed rank selection explicitly. The original task still requires economically defensible productivity envelopes or instruments, social externality measurement and a production/compensation model that survives the governance intervention. None is furnished by high pay alone.

Benjamin Moll, \emph{Lecture 11: Good Research Topics and Open Questions}, Distributional Macroeconomics lecture slides, primary PDF inspected, especially the top-income discussion around PDF page 7:
\url{https://benjaminmoll.com/wp-content/uploads/2019/07/Lecture11_2149.pdf}.

\end{document}
